\section{연구 결과}

연구 문제의 순서에 맞추어 결과를 제시하세요. 지구과학 연구에서는 자료의 시공간 범위, 불확실성, 분석 조건을 함께 설명하고, 지구과학교육 연구에서는 학습 결과 또는 교육적 의미가 드러나도록 서술합니다.

그림과 표는 학회 규정에 맞추어 간결하고 독립적으로 이해될 수 있게 작성하세요. 지구과학 자료를 제시할 때는 축 이름, 단위, 자료 출처, 분석 기간을 빠뜨리지 않는 것이 좋습니다.

한국지구과학회지는 최종 조판에서 2단 편집을 사용하므로, 표나 그림을 만들 때 처음부터 폭을 1.0 textwidth 또는 0.5 textwidth 중 하나로 정해 두는 것이 좋습니다. 지도, 시계열, 복합 패널, 넓은 표처럼 두 단을 모두 써야 읽히는 자료는 1.0 textwidth를 기준으로 만들고, 한 단 안에서 충분히 읽히는 단순 그래프나 작은 표는 0.5 textwidth를 기준으로 만듭니다. 이렇게 정해 두면 최종 조판 단계에서 축 라벨, 단위, 범례, 지도 주석이 과도하게 축소되는 문제를 줄일 수 있습니다.

\begin{ManuscriptWideTable}
\caption{한국지구과학회지 원고용 표 예시}
\label{tab:example}
\centering
\begin{tabular}{ll}
\toprule
구분 & 내용 \\
\midrule
예시 & 자료, 분석 범주, 또는 연구 문제별 결과를 입력하세요. \\
\bottomrule
\end{tabular}
\end{ManuscriptWideTable}
